% Generated draft: requires mathematical and source review.
\documentclass[11pt]{article}
\usepackage[margin=25mm]{geometry}
\usepackage{fontspec}
\setmainfont{Latin Modern Roman}
\usepackage{amsmath,amssymb,mathtools,unicode-math}
\setmathfont{Latin Modern Math}
\usepackage{enumitem,needspace,xurl,hyperref}
\setlength{\parindent}{0pt}
\setlength{\parskip}{5pt}
\setlength{\emergencystretch}{3em}
\newcommand{\N}{\mathbb{N}}
\newcommand{\Z}{\mathbb{Z}}
\newcommand{\Q}{\mathbb{Q}}
\newcommand{\R}{\mathbb{R}}
\newcommand{\C}{\mathbb{C}}
\newcommand{\F}{\mathbb{F}}
\newcommand{\A}{\mathbb{A}}
\newcommand{\PP}{\mathbb{P}}
\newcommand{\E}{\mathbb{E}}
\newcommand{\eps}{\varepsilon}
\newcommand{\dd}{\,\mathrm{d}}
\begin{document}
\section*{UnsolvedMath ID: 800883}
\textbf{Source problem code:} catalogue entry
\subsection{Definitions and mathematical statement}
Let \((X,\omega)\) be a compact K\"ahler manifold of complex dimension \(n\geq 1\), and put \(V=\int_X\omega^n\). Fix the convention \(dd^c=\frac{i}{\pi}\partial\bar\partial\). Define
\[
\operatorname{PSH}(X,\omega)=\{\varphi\in L^1(X):\varphi\text{ is upper semicontinuous and }\omega+dd^c\varphi\geq 0\},
\]
where positivity is understood in the sense of currents. For bounded \(\omega\)-plurisubharmonic functions, the measure \((\omega+dd^c\varphi)^n\) is understood in the Bedford--Taylor sense.

Let \(\mu\) be a Borel probability measure on \(X\). Interpret the stated local domination hypothesis as follows: for every \(x\in X\), there exist a coordinate neighborhood \(U\) of \(x\), a bounded plurisubharmonic function \(u\) on \(U\), and a constant \(C>0\) such that
\[
\mu|_U\leq C(dd^c u)^n
\]
as measures on \(U\). The constant can equivalently be absorbed into \(u\) by replacing it with \(C^{1/n}u\). This spells out a standard interpretation of local domination; the supplied excerpt does not define the phrase explicitly.

With the normalized convention
\[
MA_\omega(\varphi):=\frac{1}{V}(\omega+dd^c\varphi)^n,
\]
the question attributed to Ko\l{}odziej is: for every such \((X,\omega)\) and \(\mu\), does there exist
\[
\varphi\in\operatorname{PSH}(X,\omega)\cap L^\infty(X)
\qquad\text{such that}\qquad
\mu=MA_\omega(\varphi)?
\]
The supplied source writes \(MA(\varphi)\) without specifying its normalization. If it denotes the unnormalized measure \((\omega+dd^c\varphi)^n\), the probability hypothesis requires \(V=1\); that volume assumption is absent from the supplied excerpt. Thus the displayed normalized formulation is an explicitly stated interpretation, and the intended source convention remains to be checked.
\subsection{Short English statement}
The question attributed to Ko\l{}odziej asks whether local domination of a probability measure by Monge--Amp\`ere measures of bounded plurisubharmonic functions guarantees a global bounded \(\omega\)-plurisubharmonic potential. The normalization of the global Monge--Amp\`ere operator is not specified in the supplied excerpt. The registry labels this question open; that label is supplied source data, not an independently verified status assessment.
\subsection{Sources}
\begin{itemize}
\item Local UnsolvedMath registry, permanent ID \(800883\), exact source problem code \texttt{catalogue entry}, title: \emph{Open Problems in Pluripotential Theory --- Question 14}. The supplied record labels the problem open.
\item Dinew, Guedj, and Zeriahi, \emph{Open Problems in Pluripotential Theory} (2015), Question 14, attributed there to Ko\l{}odziej. Supplied source URL: \url{https://arxiv.org/abs/1511.00705}.
\item The supplied provenance records access on \texttt{2026-10-16} and extraction by \texttt{arxiv\_latex}. These are supplied provenance claims; no source visit or independent verification is asserted here.
\end{itemize}
\end{document}
